\chapter{The Story of Telecommunication}
\label{ch:history}

\begin{chapterintro}
Communication engineering is a young discipline built on old problems. This chapter
follows the problem of sending messages at a distance from signal fires to 5G, and at each
step extracts the technical idea that survives in today's equipment: the code that
compresses text, the cable that smears pulses, the amplifier that made long distance
possible, the carrier that let many conversations share one wire, the receiver
architecture inside every radio, the theorems that define what is possible, and the
engineering choices that shaped satellites, cellular networks and the Internet. Read it as
a map of the rest of the book.
\end{chapterintro}

\section{Why an engineer should know the history}

It is possible to design a modem without knowing who Harry Nyquist was, just as it is
possible to write software without knowing why a byte has eight bits. But communication
systems are unusually shaped by their past. The 3.4\,kHz bandwidth of a telephone
channel, the 8\,kHz sampling rate of digital voice, the 1.544\,Mb/s of a T1 line, the
10.7\,MHz intermediate frequency of an FM radio, the 15\,kHz subcarrier spacing of LTE and
5G, and the 30\,kHz channels of the first cellular systems are all decisions made by
engineers under particular constraints, and they propagate forward because new systems
must coexist with old ones. When you understand why a number was chosen, you know when
you may safely change it.

History also teaches humility. The transatlantic cable of 1858 failed because its
engineers did not understand the physics of a long line. The early theorists of frequency
modulation concluded that it had no advantage, and they were right about the question
they asked and wrong about the one that mattered. Code-division multiple access was
declared impractical by experts shortly before it became the basis of third-generation
cellular. Low-density parity-check codes were invented in 1960, forgotten for thirty-five
years, and are now in every 5G phone. Engineering is a conversation across generations,
and it is worth hearing the earlier voices.

\section{Before electricity}

The oldest long-distance communication systems were visual. Greek historians describe
chains of signal fires; the Roman army used towers; and the Great Wall of China had
beacon towers that could relay a warning hundreds of kilometres in hours. These systems
had enormous speed---the message moved at the speed of light between towers---but
almost no \emph{capacity}. A fire is either lit or not, so the message had to be agreed in
advance: ``the enemy is coming.'' In modern language, the alphabet had two symbols and the
system sent one bit per event.

\subsection{The optical telegraph}

The first system that could send arbitrary messages over long distances was the optical
telegraph of Claude Chappe. Beginning with the Paris--Lille line in 1794, the French
government built a network of towers topped by articulated arms whose positions could be
read by telescope from the next tower, some 10 to 15\,km away. Chappe coined the word
\emph{t\'el\'egraphe}, ``far writer.'' The arms could take 98 distinguishable positions, of which a
subset was used for control and the rest indexed a codebook of words and phrases: a
message was a sequence of page and line numbers. At its height the French network had
more than 500 stations, and a message could cross the country in a few hours.

The Chappe telegraph already contained ideas that recur throughout this book. It used a
finite \emph{alphabet} of signal states chosen to be distinguishable in the presence of
``noise''---haze, fatigue and bad telescopes. It used a \emph{codebook} to map frequent
messages into short symbol sequences. It had a \emph{protocol}: control signals for
``start'', ``error'' and ``repeat''. And it had a \emph{capacity limit}, set by how quickly an
operator could move the arms and how reliably the next operator could read them, of
roughly one symbol every 20 to 30 seconds.

\section{The electric telegraph}

\subsection{Cooke, Wheatstone, Morse and Vail}

Electricity offered a signal that moved along a wire faster than any horse and did not
care about fog. William Cooke and Charles Wheatstone patented a needle telegraph in
Britain in 1837; it used several needles that pointed at letters on a board, so an untrained
operator could read it. In the United States, Samuel Morse and his collaborator Alfred Vail
chose a different path: a single circuit, switched on and off by a key, whose pattern of
short and long pulses encoded letters. On 24 May 1844 Morse sent ``What hath God
wrought'' from Washington to Baltimore, and within a decade telegraph wires followed the
railways across the continent.

The Morse code is a remarkable piece of engineering. Its designers wanted to minimise the
time needed to send typical English text, and they realised that the answer was to give
short codes to frequent letters. According to the traditional account, Vail estimated letter
frequencies by counting the type in a printer's type case: the printer kept 12\,000 E's
but only 200 Z's. The code gives E a single dot and T a single dash, while Q, Z and J
receive four symbols. This is a \emph{variable-length source code}, a century before David
Huffman proved (in 1952) how to construct the optimal one, and Morse code is within about
15\% of the efficiency of a Huffman code for English letters. We will return to this in
Chapter~\ref{ch:sourcecoding}.

The timing rules of International Morse code are strict: a dot lasts one unit, a dash three,
the gap between symbols one, between letters three and between words seven.
Figure~\ref{fig:ch01_morse} shows the keyed envelope of ``SOS'' and something the early
radio operators discovered the hard way: the spectrum of the keyed signal. Switching a
carrier on and off instantaneously produces sharp edges, and sharp edges have wide
spectra. A radio telegraph with hard keying spattered interference---``key clicks''---across
neighbouring channels. The cure was to shape the edges of each pulse, and the lower
curve shows that raised-cosine edges lower the spectral sidelobes by tens of decibels.
Pulse shaping to control bandwidth is exactly the subject of Chapter~\ref{ch:baseband},
where it becomes the root-raised-cosine filter inside every modern modem.

\mdcfig{ch01_morse}{Top: the on-off keyed envelope of ``SOS''. Bottom: the power spectrum of
repeated Morse traffic with instantaneous (hard) keying and with raised-cosine shaped
edges. Shaping the pulses reduces the out-of-band spectrum by more than 40\,dB, which is why
operators added ``click filters'' to their keys.}

\subsection{Multiplexing: sharing one wire}

Telegraph wires were expensive, and a single operator could send only 30 or 40 words per
minute. The first attack on this bottleneck was \emph{duplex} and \emph{quadruplex}
operation: Thomas Edison's quadruplex of 1874 sent two messages in each direction at
once by using both the polarity and the magnitude of the line current as independent
signals---an early form of what we now call multilevel signalling and full-duplex operation.

The second attack was time-division multiplexing. \'Emile Baudot's printing telegraph of
1874 used a rotating distributor that connected several operators to the line in turn, and
each operator entered characters using five keys. Every character was therefore a
fixed-length five-bit word, the Baudot code, the ancestor of ASCII. The unit of signalling
rate, the \emph{baud}, honours him: one baud is one symbol per second. Note the
distinction, which we will use constantly: the baud rate counts symbols, while the bit rate
counts bits, and the two differ whenever a symbol carries more than one bit.

\begin{keyidea}[Three ways to share a channel]
Every multiplexing scheme ever built divides the channel's resources along some
dimension so that the signals do not interfere. Baudot divided \emph{time}. Carrier
telephony (below) divided \emph{frequency}. Code-division systems, a century later, divided
the signal space using nearly orthogonal \emph{codes}, and multi-antenna systems divide
\emph{space}. Chapter~\ref{ch:multipleaccess} treats all four as one idea: making the
signals of different users approximately orthogonal.
\end{keyidea}

\subsection{The Atlantic cable and the physics of a long line}

The most dramatic engineering story of the nineteenth century is the laying of the
transatlantic telegraph cable. The first cable, completed in August 1858, carried Queen
Victoria's 98-word message to President Buchanan---which took about sixteen hours to
transmit. Within weeks the cable failed entirely. The company's electrician, Wildman
Whitehouse, had tried to overcome the weak, sluggish signals by applying about 2000 volts,
and the insulation broke down.

The man who understood what was really happening was William Thomson, later Lord
Kelvin. In 1855 he had analysed a long submarine cable as a distributed resistance $R$ and
capacitance $C$ per unit length. Neglecting inductance and leakage, the voltage obeys a
diffusion equation,
\begin{equation}
\frac{\partial v}{\partial t} = \frac{1}{RC}\,\frac{\partial^2 v}{\partial x^2},
\label{eq:ch01:diffusion}
\end{equation}
whose response at distance $\ell$ to a voltage step applied at $t=0$ is
\begin{equation}
v(\ell,t) = \operatorname{erfc}\!\left(\frac{\ell}{2}\sqrt{\frac{RC}{t}}\right).
\end{equation}
The time for the received voltage to reach a given fraction of its final value grows as
$RC\ell^2$: double the length of the cable and you must signal four times more slowly. This
is Thomson's ``law of squares''. Figure~\ref{fig:ch01_cable} shows the consequence. When
dots are sent slowly, each has time to rise and fall; when they are sent faster, the slow
tails of earlier pulses pile up on later ones and the pattern is lost. The name for this is
\term{intersymbol interference} (ISI), and it is the central problem of
Chapters~\ref{ch:baseband} and~\ref{ch:equalization}.

\mdcfig{ch01_cable}{Kelvin's model of a long submarine cable. Left: the response to a step
of voltage at three cable lengths; the delay scales with the square of the length. Right:
a train of pulses through the same cable at two keying speeds. At the faster rate each
pulse is smeared into its neighbours---intersymbol interference.}

Thomson's remedies were engineering, not force. He invented the sensitive mirror
galvanometer so that small currents could be read, and he argued for a thicker conductor
and better insulation to lower $R$ and $C$. The 1866 cable, laid by the steamship
\emph{Great Eastern}, embodied his design and worked for years, carrying about eight words
per minute. Thomson was knighted for it.

\begin{history}[Heaviside and the distortionless line]
Oliver Heaviside, a self-taught telegraph engineer, extended Thomson's analysis in the
1880s by including the inductance $L$ and leakage conductance $G$ of the line. He found
that when $L/R = C/G$, all frequencies travel at the same speed and are attenuated
equally: the line becomes \emph{distortionless}. Real lines had far too little inductance, so
Heaviside proposed adding it. His suggestion was ignored for years, partly because of his
combative personality, until Michael Pupin and George Campbell independently patented
\emph{loading coils} around 1900. Loaded telephone cables, with coils every 1.8\,km, made
long-distance telephony practical---and a century later, those same coils had to be
removed from the copper loop before DSL could work, because they turn the line into a
low-pass filter that blocks everything above the voice band (Chapter~\ref{ch:wireline}).
\end{history}

\section{The telephone and the network}

\subsection{Voice over a wire}

Alexander Graham Bell's telephone patent was filed on 14 February 1876, the same day
Elisha Gray filed a caveat on a similar idea. Bell's transmitter converted the pressure
variations of speech into a proportional current: an \emph{analog} signal. The telegraph
had sent discrete symbols; the telephone sent a continuous waveform, and it demanded a
new kind of engineering. A telegraph signal only had to be detectable; a telephone
signal had to be reproduced faithfully, which meant controlling the frequency response,
noise and distortion of every component in the path.

Experiments with listeners established that intelligible speech needs roughly the band
from 300 to 3400\,Hz. That number---about 3.1\,kHz of useful bandwidth inside a 4\,kHz
channel allocation---was fixed by the 1920s and still defines the digital telephone
network: the 8\,kHz sampling rate of pulse-code modulation (Chapter~\ref{ch:pcm}) is twice
the 4\,kHz channel width.

\subsection{Switching and the first automation}

Early telephones were connected through manual switchboards. The first automatic switch
was invented by Almon Strowger, an undertaker in Kansas City who, according to the story,
suspected that the local operator was diverting his customers to a competitor. His
step-by-step switch of 1891 moved a wiper in response to the pulses of a rotary dial:
dialling ``3'' produced three pulses that stepped the switch three positions. The
telephone network became the largest machine on Earth, and its growth drove the
development of switching theory, traffic theory (A.~K.~Erlang's queueing formulas, which
still size cellular networks; Chapter~\ref{ch:multipleaccess}) and ultimately the digital
computer.

\subsection{Amplification and the long-distance problem}

Loading coils extended the reach of a telephone call to perhaps 1500\,km. Beyond that
the signal simply became too weak. The solution was the \emph{repeater}, an amplifier
inserted into the line, and the device that made it possible was the vacuum tube. Lee de
Forest's Audion triode (1906) could amplify, but only erratically; engineers at AT\&T led by
Harold Arnold turned it into a reliable high-vacuum amplifier, and on 25 January 1915 the
first transcontinental telephone call connected New York and San Francisco through a
chain of tube repeaters.

Amplifiers created a new problem. Every repeater added a little distortion, and a
transcontinental circuit had dozens of them. In August 1927 Harold Black, a young Bell
Labs engineer, sketched on a page of the \emph{New York Times} during his ferry commute the
idea of the \term{negative feedback amplifier}: sacrifice most of the amplifier's gain by
feeding an inverted fraction of the output back to the input, and in exchange the overall
gain depends almost entirely on the stable passive feedback network rather than on the
drifting tube. Negative feedback is now so universal that it is hard to imagine
electronics without it, and the stability analysis it required---by Harry Nyquist in 1932
and Hendrik Bode---founded control theory.

\subsection{Carrier telephony: frequency-division multiplexing}

Once a wire could carry frequencies far above the voice band, it was natural to put many
conversations on it at once, each shifted to its own frequency slot. This is
\term{frequency-division multiplexing} (FDM). Its key enabler was \emph{single-sideband}
modulation, patented by John Carson of AT\&T in 1915: by transmitting only one of the two
sidebands of an amplitude-modulated carrier, each voice channel occupied only its own
4\,kHz. The first carrier system went into service in 1918. By the 1940s the Bell
System's L-carrier on coaxial cable stacked 600 voice channels, and by the 1970s the L5
system carried 10\,800 voice channels on each coaxial tube, organised into a hierarchy of
groups (12 channels), supergroups (60), mastergroups (600) and beyond.
Chapter~\ref{ch:analog} builds this hierarchy from first principles.

\section{Wireless}

\subsection{From Maxwell to Marconi}

In 1865 James Clerk Maxwell published the equations that unified electricity, magnetism
and light and predicted electromagnetic waves travelling at the speed of light. Heinrich
Hertz generated and detected such waves in his laboratory between 1886 and 1888, using a
spark gap as a transmitter and a small loop with a gap as a receiver. He thought the
discovery had no practical use.

Guglielmo Marconi thought otherwise. Beginning in 1895, he improved the range of
spark-gap systems by raising the antennas and grounding them, and on 12 December 1901 he
reported receiving the letter S in Morse code at Signal Hill, Newfoundland, sent from
Poldhu in Cornwall---across 3500\,km of ocean and, unknown to him, bent around the Earth's
curvature by the ionosphere. Wireless telegraphy became essential at sea. The sinking of
the \emph{Titanic} in 1912, when a nearby ship's radio operator had gone to bed, led directly
to international rules requiring continuous radio watch and to the first licensing of
radio stations.

Spark transmitters were extremely wasteful. A spark produces a damped oscillation whose
spectrum spreads over a wide band, so only a few stations could operate in one area. The
next generation used continuous waves from high-frequency alternators and arc
converters, and eventually from vacuum-tube oscillators. A continuous-wave transmitter
concentrates its energy at one frequency, which allows receivers to be selective and many
stations to coexist: the start of spectrum management.

\subsection{Voice by radio, and Armstrong's three inventions}

Reginald Fessenden, a Canadian engineer, understood that a continuous wave could be
\emph{modulated} by voice, and he reportedly broadcast speech and music on Christmas Eve 1906
from Brant Rock, Massachusetts. Amplitude modulation (AM) became the basis of
broadcasting; the first scheduled commercial broadcast is usually credited to KDKA in
Pittsburgh, which reported the results of the Harding--Cox presidential election on
2 November 1920.

The receivers of the day were poor, and three inventions by Edwin Howard Armstrong fixed
them. The first, the \emph{regenerative receiver} (1912), fed part of a tube amplifier's output
back to its input with positive feedback, giving enormous gain from one tube. The second,
the \term{superheterodyne receiver} (1918), mixed every incoming signal down to a fixed
\emph{intermediate frequency} (IF), where selective, high-gain amplifiers could be built once
and optimised. Nearly every radio built since---including the software-defined radio in
your lab, which is a superheterodyne whose IF happens to be zero---follows this
architecture. The third, wideband \emph{frequency modulation} (FM, patented 1933),
we return to shortly.

\begin{keyidea}[The superheterodyne principle]
It is hard to build a filter that is both very selective and tunable over a wide range. It
is easy to build a very selective filter at one fixed frequency. The superhet therefore
moves the signal to the filter rather than moving the filter to the signal, using a mixer
and a tunable local oscillator. Chapter~\ref{ch:analog} derives the image-frequency problem
that this creates; Chapter~\ref{ch:sdr} shows how modern zero-IF and low-IF receivers
solve it digitally.
\end{keyidea}

\subsection{FM: the answer to a different question}

In 1922 John Carson analysed frequency modulation as a way of saving bandwidth and
showed, correctly, that it could not: an FM signal always occupies at least as much
bandwidth as the equivalent AM signal. He concluded that FM ``inherently distorts without
any compensating advantages.'' Armstrong asked a different question: what if we
deliberately use \emph{more} bandwidth? His 1936 paper showed that wideband FM trades
bandwidth for signal-to-noise ratio: a broadcast FM signal that uses about 180\,kHz to carry
15\,kHz of audio is dramatically quieter than AM. This was the first practical
demonstration of a principle that Shannon would make exact twelve years later: bandwidth
and signal power are interchangeable currencies. FM broadcasting, the audio of analog
television, two-way police and military radio, and the first cellular systems were all
built on it.

\section{The theorists}

While engineers built networks, a handful of theorists at Bell Labs and elsewhere worked
out what the networks could possibly do. Their results are the foundation of this book.

\begin{description}[leftmargin=0pt,itemsep=4pt]
\item[Harry Nyquist] showed in 1924 and 1928 that a channel of bandwidth $B$ can carry at
most $2B$ independent pulses per second without intersymbol interference, and derived the
condition a pulse must satisfy to achieve it---the \emph{Nyquist criterion} of
Chapter~\ref{ch:baseband}. With John B.~Johnson he also explained thermal noise in 1928:
every resistor at temperature $T$ generates noise with power $kTB$, the floor under every
radio receiver (Chapter~\ref{ch:noise}).
\item[Ralph Hartley] proposed in 1928 that information be measured as the logarithm of
the number of distinguishable messages, so that a channel sending $2B$ symbols per second
from an alphabet of $M$ levels carries $2B\log_2 M$ bits per second.
\item[Vladimir Kotelnikov] in the Soviet Union (1933), following E.~T.~Whittaker (1915),
stated the sampling theorem: a signal band-limited to $B$ is completely determined by
samples taken $2B$ times per second.
\item[Norbert Wiener] and Andrey Kolmogorov developed, during the Second World War, the
theory of optimal linear filtering of noisy signals, the ancestor of the MMSE equalizer.
\item[Stephen O.~Rice] of Bell Labs published in 1944--45 the mathematical analysis of
random noise that gave the Rayleigh and Rician distributions their place in fading
channel theory (Chapter~\ref{ch:channels}).
\end{description}

\subsection{Claude Shannon}

In 1948 Claude Shannon, then 32 and working at Bell Labs, published ``A Mathematical
Theory of Communication.'' It is one of the rare papers that created a field fully formed.
Shannon modelled communication with the block diagram of Figure~\ref{fig:ch01:shannon},
defined information as a measurable quantity---the entropy of the source, in bits---and
proved two theorems that define the limits of the field. The \emph{source coding theorem}
says that a source can be compressed to its entropy and no further. The \emph{channel
coding theorem} says that every channel has a capacity $C$: at any rate below $C$ it is
possible to communicate with arbitrarily small error probability, and at any rate above it
it is not. For the band-limited Gaussian channel,
\begin{equation}
C = B\log_2\!\left(1 + \frac{S}{N}\right)\quad\text{bits per second.}
\label{eq:ch01:shannon}
\end{equation}

\begin{figure}[htbp]\centering
\begin{tikzpicture}[node distance=0.55cm and 0.55cm,
  blk/.style={draw=navy,fill=navy!6,thick,minimum height=1.0cm,minimum width=1.75cm,align=center,font=\sffamily\small},
  arr/.style={-{Stealth[length=2.5mm]},thick,navy}]
\node[blk] (src) {Information\\source};
\node[blk,right=of src] (tx) {Transmitter};
\node[draw=accent,circle,thick,right=of tx,minimum size=0.7cm,font=\sffamily\small] (ch) {$+$};
\node[blk,right=of ch] (rx) {Receiver};
\node[blk,right=of rx] (dst) {Destination};
\node[blk,above=0.55cm of ch,fill=accent!8,draw=accent] (ns) {Noise\\source};
\draw[arr] (src) -- node[above,font=\scriptsize]{message} (tx);
\draw[arr] (tx) -- node[above,font=\scriptsize]{signal} (ch);
\draw[arr] (ch) -- node[above,font=\scriptsize]{received} (rx);
\draw[arr] (rx) -- node[above,font=\scriptsize]{message} (dst);
\draw[arr,accent] (ns) -- (ch);
\end{tikzpicture}
\caption{Shannon's model of a communication system (redrawn after his 1948 paper). Every
system in this book---telegraph, telephone, radio, satellite, fibre, cellular---fits this diagram.}
\label{fig:ch01:shannon}
\end{figure}

Shannon's proof was not constructive: it showed that good codes exist by averaging over
randomly chosen codes, and it said nothing about how to build or decode one. The rest of
the history of coding theory, from Richard Hamming's single-error-correcting codes of 1950
to polar codes in 2009, is the story of the search for practical codes that approach
Shannon's limit. Chapters~\ref{ch:infotheory}--\ref{ch:moderncodes} tell it.

\begin{keyidea}[Shannon's lesson for the working engineer]
Before designing anything, compute the capacity of the channel you have. It tells you
whether your goal is impossible, easy or somewhere in between, and it tells you whether
to spend your effort on more power, more bandwidth, more antennas or a better code. Every
link budget in this book ends with that comparison.
\end{keyidea}

\section{War, secrecy and the birth of digital communication}

The Second World War accelerated communication technology as much as any peacetime
decade. Three developments stand out.

\emph{Pulse-code modulation}, invented by Alec Reeves in 1937, represented a voice signal
as a sequence of numbers. At first it seemed absurdly wasteful of bandwidth, but it had one
decisive advantage: a digital signal can be \emph{regenerated}. An analog signal accumulates
noise at every repeater; a digital signal can be detected and re-transmitted perfectly
clean, so noise does not accumulate. The first operational use of digital voice was
SIGSALY, the secure telephone system that connected Roosevelt and Churchill from 1943,
which combined a vocoder, quantized samples and one-time-pad encryption from phonograph
records.

\emph{Frequency hopping} was patented in 1942 by the actress Hedy Lamarr and the composer
George Antheil, who proposed synchronising the frequency changes of a torpedo's radio
link using piano-roll mechanisms, so that an enemy could not jam it. The patent was not used
during the war, but the principle of spreading a signal over a wide band to make it
resistant to jamming and interception---\term{spread spectrum}---became central to military
communications, GPS, Bluetooth and code-division cellular systems
(Chapter~\ref{ch:spreadspectrum}).

\emph{Radar and microwaves.} The cavity magnetron, brought from Britain to the United
States in 1940, made high-power microwave transmitters practical, and the MIT Radiation
Laboratory developed the microwave engineering---waveguides, mixers, low-noise
receivers, the theory of detection in noise---that post-war microwave relay links and
satellite communication were built on. The \emph{matched filter}, which maximises the
signal-to-noise ratio of a known pulse in noise, was derived by D.~O.~North for radar in
1943; it is the optimal receiver front end of Chapter~\ref{ch:baseband}.

\section{Transistors, satellites and fibre}

\subsection{The solid-state revolution}

In December 1947 John Bardeen, Walter Brattain and William Shockley at Bell Labs
demonstrated the point-contact transistor. The integrated circuit followed in 1958--59
(Jack Kilby at Texas Instruments, Robert Noyce at Fairchild), and with it the
exponential scaling that would eventually make it cheaper to process a radio signal
digitally than with analog components. The telephone network's first large digital
transmission system, the T1 carrier of 1962, sent 24 PCM voice channels over two twisted
pairs at 1.544\,Mb/s, with regenerative repeaters every 1.8\,km---exactly where the loading
coils had been.

\subsection{Communication satellites}

In a 1945 article in \emph{Wireless World}, the science-fiction writer and radar engineer
Arthur C.~Clarke observed that a satellite at an altitude of about 35\,786\,km above the
equator orbits once per sidereal day and therefore appears fixed in the sky; three such
satellites could cover almost the whole Earth. Twelve years later the Soviet Union launched
Sputnik~1 (4 October 1957), whose radio beacons on 20 and 40\,MHz were tracked by
amateurs worldwide---and whose Doppler shift, measured at Johns Hopkins, led directly to
the Transit satellite navigation system and eventually to GPS.

Progress was then rapid. NASA's Echo~1 (1960) was a passive 30\,m metallised balloon that
simply reflected signals. AT\&T's Telstar~1 (1962) carried the first live transatlantic
television, but in a low orbit it was visible to its ground stations for only minutes at a
time. Syncom~2 (1963) reached geosynchronous orbit, and Syncom~3 (1964), in true
geostationary orbit, relayed television of the Tokyo Olympics to the United States.
Intelsat~I, ``Early Bird'', began commercial transatlantic service in 1965 with 240 telephone
circuits. Chapter~\ref{ch:satellite} develops orbital mechanics, satellite link budgets and
modern systems from DVB-S2 broadcasting to the low-Earth-orbit constellations that, sixty
years after Telstar, have returned to Telstar's orbit in their thousands.

\subsection{Light}

In 1966 Charles Kao and George Hockham argued that the high loss of glass fibre was due to
impurities rather than to anything fundamental, and that purified silica could reach
losses below 20\,dB/km, making optical fibre competitive with coaxial cable. Corning
achieved that figure in 1970; modern fibre loses about 0.17--0.2\,dB/km at 1550\,nm. The
erbium-doped fibre amplifier (1987) amplified light directly, without converting it to
electricity, and allowed many wavelengths to be amplified at once. The first transatlantic
fibre cable, TAT-8 (1988), carried 280\,Mb/s; a modern submarine cable carries hundreds of
terabits per second using coherent detection, the same QAM and DSP you will study in
Chapters~\ref{ch:modulation}--\ref{ch:equalization}, applied to light
(Chapter~\ref{ch:wireline}). Kao received the Nobel Prize in Physics in 2009.

\section{Networks of computers}

The ARPANET connected its first two nodes in 1969, and Robert Metcalfe's Ethernet (1973)
connected computers in a building by letting them share a coaxial cable with a
listen-before-talk protocol derived from the University of Hawaii's ALOHAnet radio network.
The data modems that carried computer traffic over the telephone network are a history of
physical-layer engineering in miniature: from 300\,b/s frequency-shift keying in the 1960s
to 9600\,b/s with 32-point QAM and trellis-coded modulation (Gottfried Ungerboeck's 1982
invention) in V.32, to 33.6\,kb/s in V.34---within a few decibels of the Shannon capacity of a
telephone channel. Figure~\ref{fig:ch01_rates} puts these numbers in perspective.

\mdcfig{ch01_rates}{Approximate data rates of landmark systems. Each family advanced by
roughly a factor of ten every six to ten years once it was digital; the values are
order-of-magnitude illustrations of typical or peak rates rather than precise
specifications.}

\section{Cellular radio}

The cellular concept was described in a 1947 Bell Labs memorandum by Douglas Ring and
Rae Young: instead of one powerful transmitter covering a city, use many low-power base
stations, each covering a small \emph{cell}, and reuse the same frequencies in cells far
enough apart. Capacity then grows with the number of cells rather than being fixed by the
spectrum. The idea waited three decades for the technology---microprocessors to manage
handoff between cells, and spectrum allocations to use.

\begin{table}[htbp]\centering\small
\caption{Five generations of cellular technology.}\label{tab:ch01:gens}
\begin{tabularx}{\linewidth}{@{}l l l X@{}}
\toprule
Generation & Launch & Key systems & Defining physical-layer ideas\\
\midrule
1G & 1979--83 & NTT (Tokyo), AMPS, NMT, TACS & Analog FM voice, 25--30\,kHz FDMA channels, handoff\\
2G & 1991--95 & GSM, IS-95 (cdmaOne), IS-136 & Digital speech codecs, TDMA with GMSK (GSM), CDMA with power control and soft handoff (IS-95), SMS\\
3G & 2001 & UMTS/WCDMA, CDMA2000 & Wideband CDMA, turbo codes, fast power control, HSPA packet data\\
4G & 2009 & LTE, LTE-Advanced & OFDMA downlink, SC-FDMA uplink, MIMO, all-IP, carrier aggregation\\
5G & 2019 & NR & Flexible OFDM numerology, LDPC and polar codes, massive MIMO, mmWave beamforming\\
6G & $\approx$2030 & (Release 21 onward) & Evolved OFDM, sensing, AI-native design, non-terrestrial integration\\
\bottomrule
\end{tabularx}
\end{table}

The first commercial cellular network opened in Tokyo in 1979; in the United States,
the Advanced Mobile Phone System (AMPS) launched in Chicago in October 1983, six and a half
years after Martin Cooper of Motorola had made the first handheld cellular call on a New
York street in April 1973. Table~\ref{tab:ch01:gens} summarises what followed. Each
generation brought the newest physical-layer ideas of its time into billions of devices:
digital modulation and speech coding in 2G, code-division multiple access and turbo
codes in 3G, OFDM and MIMO in 4G, LDPC and polar codes and massive antenna arrays in 5G.
Chapter~\ref{ch:cellular} tells the full story, including the ``CDMA wars'' of the early
1990s, when a small company named Qualcomm---founded by Irwin Jacobs and Andrew Viterbi,
the inventor of the Viterbi algorithm---persuaded a sceptical industry that CDMA could work.

\section{Unlicensed radio and Wi-Fi}

In 1985 the US Federal Communications Commission opened the industrial, scientific and
medical (ISM) bands at 902\,MHz, 2.4\,GHz and 5.8\,GHz to unlicensed spread-spectrum
communication. Nobody expected much of it: 2.4\,GHz is where microwave ovens heat water.
The IEEE 802.11 standard followed in 1997 at 1--2\,Mb/s, 802.11b reached 11\,Mb/s in 1999,
and the industry adopted the brand name Wi-Fi. Today more Internet traffic reaches its final
destination over Wi-Fi than over any other medium, and the bands that were once a spectral
junkyard are among the most valuable in the world. The lesson for the spectrum engineer is
that flexible rules and cheap hardware can create more value than exclusive licences.
Chapter~\ref{ch:wifi} follows 802.11 from 1\,Mb/s to Wi-Fi~7's 46\,Gb/s peak.

\section{The coding renaissance and the MIMO revolution}

Between 1993 and 2009, coding theory went from ``several decibels from capacity'' to
``within a fraction of a decibel''. In 1993 Claude Berrou, Alain Glavieux and Punya
Thitimajshima presented \emph{turbo codes}, which used two simple convolutional codes and
an iterative decoder to come within 0.7\,dB of the Shannon limit; many experts initially
assumed the results were wrong. In 1996 David MacKay and Radford Neal rediscovered Robert
Gallager's 1960 \emph{low-density parity-check} codes and showed they were as good. And in
2009 Erdal Ar{\i}kan introduced \emph{polar codes}, the first codes proven to achieve
capacity with low-complexity encoding and decoding. 5G uses both LDPC and polar codes.

At the same time, Gerard Foschini at Bell Labs (1996) and Emre Telatar (1995) showed that a
channel with multiple antennas at both ends has a capacity that grows linearly with the
number of antennas---an entirely new resource, obtained from the same bandwidth and power
by exploiting the multipath scattering that radio engineers had always regarded as a
nuisance. Siavash Alamouti's two-antenna space-time code (1998) was simple enough to put
into every handset. Multiple-input multiple-output (MIMO) transmission is now in every
Wi-Fi access point and cellular base station (Chapter~\ref{ch:mimo}).

\section{The software-defined radio}

For most of the twentieth century a radio \emph{was} its hardware: its filters, oscillators
and detectors fixed what it could do. As digital processing grew cheaper, it became
possible to digitise the signal close to the antenna and do everything else in software.
Joseph Mitola coined the term \emph{software radio} in 1992. GNU Radio, an open-source
framework for building radios in software, began in 2001, and the Universal Software Radio
Peripheral (USRP) from Ettus Research gave students and researchers affordable hardware for
it. Today the ``software radio'' idea is simply how radios are built: a 5G base station is
largely software running on general-purpose or specialised processors, and the Ettus B200
on your bench is architecturally a scaled-down cousin of a cellular radio unit
(Chapter~\ref{ch:sdr}).

\section{Lessons from two centuries}

Looking back over this history, a few patterns recur so often that they amount to
engineering principles.

\begin{enumerate}
\item \textbf{Every channel is eventually limited by noise and dispersion.} Kelvin's cable,
Nyquist's telegraph and a 5G link differ in scale, not in kind. Intersymbol interference and
noise are the two enemies, and every technique in this book fights one or both.
\item \textbf{Digital wins because it can be regenerated.} Once a signal is represented by
discrete symbols, errors can be prevented from accumulating and corrected by coding. This
is why every analog system has been replaced by a digital one.
\item \textbf{Bandwidth, power, antennas and complexity are interchangeable currencies.}
Armstrong traded bandwidth for SNR; Shannon quantified the exchange rate; turbo and LDPC
codes traded computation for power; MIMO traded antennas for bandwidth. Good design is
choosing the cheapest currency.
\item \textbf{Old ideas return when technology catches up.} Gallager's LDPC codes, the
cellular concept, frequency hopping and multicarrier modulation (first proposed in the
1960s) all waited decades for cheap enough electronics.
\item \textbf{Compatibility shapes everything.} The 4\,kHz voice channel, the 8\,kHz
sampling rate and the 15\,kHz subcarrier spacing persist because new systems had to live
with old ones.
\end{enumerate}

\begin{labbox}[Chapter 1 simulations]
The figure script \texttt{figscripts/ch01\_figs.py} reproduces Kelvin's cable model and the
key-click spectrum. Try changing the cable length and keying speed until the pulse pattern
becomes unreadable, and estimate the maximum ``words per minute'' of a cable twice as long.
Lab~1 (\lab{lab01\_baseband.ipynb}) then moves from telegraph pulses to the complex baseband
signals of a modern radio.
\end{labbox}

\problems
\begin{enumerate}[label=\thechapter.\arabic*]
\item Using Thomson's law of squares, if the 1866 cable (about 3000\,km) supported 8 words
per minute, what rate would a 1500\,km cable of the same construction support? What would
halving the capacitance per unit length achieve?
\item The Chappe telegraph could display 98 positions, of which about 92 carried data, at
about one position every 20\,s. Compute its information rate in bits per second using
Hartley's measure. Compare with a Morse operator sending 20 words per minute (assume 5
letters per word and 4.7 bits per letter of English text).
\item Morse code assigns E the shortest code and Q a long one. Using the timing rules in
the text, compute the average duration in dot units of a letter of English, given the
approximate letter frequencies E 12.7\%, T 9.1\%, A 8.2\%, O 7.5\%, I 7.0\%, N 6.7\%, S 6.3\%,
H 6.1\%, R 6.0\% (assign the remaining 36.4\% an average code of 3 symbols with 2 dashes).
Compare with a fixed-length code.
\item A carrier telephone system stacks 600 single-sideband voice channels, each allotted
4\,kHz. What total bandwidth does it require? If double-sideband AM had been used instead,
how many channels would fit in the same band?
\item Use Equation~\eqref{eq:ch01:shannon} to compute the capacity of a 3.1\,kHz telephone
channel with a 35\,dB signal-to-noise ratio. How close did the V.34 modem (33.6\,kb/s) come?
\item Why was the ionosphere essential to Marconi's 1901 experiment? Estimate the maximum
line-of-sight distance between two 150\,m antenna masts on a spherical Earth of radius
6371\,km.
\item \emph{(Research)} Choose one of the ``old ideas'' listed in the lessons---LDPC codes,
multicarrier modulation, frequency hopping or the cellular concept---and write a one-page
account of what technology had to change before it could be deployed.
\end{enumerate}

\furtherreading
\begin{itemize}
\item T.~Standage, \emph{The Victorian Internet}, Walker, 1998. A lively history of the
telegraph and its social impact.
\item J.~Gleick, \emph{The Information: A History, a Theory, a Flood}, Pantheon, 2011. From
talking drums to Shannon, beautifully written.
\item J.~Gertner, \emph{The Idea Factory: Bell Labs and the Great Age of American Innovation},
Penguin, 2012.
\item J.~Soni and R.~Goodman, \emph{A Mind at Play: How Claude Shannon Invented the Information
Age}, Simon \& Schuster, 2017.
\item C.~E.~Shannon, ``A Mathematical Theory of Communication,'' \emph{Bell System Technical
Journal}, 27, 1948. Still readable, and still the best introduction to its subject.
\item P.~J.~Nahin, \emph{Oliver Heaviside: The Life, Work, and Times of an Electrical Genius of
the Victorian Age}, Johns Hopkins, 2002.
\item A.~C.~Clarke, ``Extra-Terrestrial Relays,'' \emph{Wireless World}, October 1945.
\end{itemize}
